\documentclass[11pt]{article}
\usepackage[margin=1in]{geometry}
\usepackage{graphicx}
\usepackage{float}
\renewcommand{\topfraction}{0.9}\renewcommand{\bottomfraction}{0.9}\renewcommand{\textfraction}{0.08}\renewcommand{\floatpagefraction}{0.7}
\usepackage{booktabs}
\usepackage{xcolor}
\usepackage{amsmath}
\usepackage{lmodern}
\usepackage{authblk}
\usepackage[hidelinks]{hyperref}
\setlength{\parindent}{0pt}
\setlength{\parskip}{5pt}
\graphicspath{{figures/}}

% figures and tables numbered by hand-rolled counters (the caption package is not needed)
\newcounter{fignum}
\newcounter{tabnum}
\newcommand{\fig}[3]{%
  \begin{figure}[tbp]\centering
  \includegraphics[width=\linewidth]{#1}\par\vspace{2pt}
  \refstepcounter{fignum}\label{#2}%
  {\small\raggedright\textbf{Figure \thefignum.} #3\par}
  \end{figure}}
\newcommand{\tabsmall}[3]{%
  \begin{minipage}{\linewidth}\centering
  \refstepcounter{tabnum}\label{#2}%
  {\small\raggedright\textbf{Table \thetabnum.} #3\par}\vspace{3pt}
  \input{tables/#1}
  \end{minipage}\par\vspace{10pt}}
\newcommand{\tab}[3]{%
  \begin{minipage}{\linewidth}\centering
  \refstepcounter{tabnum}\label{#2}%
  {\small\raggedright\textbf{Table \thetabnum.} #3\par}\vspace{3pt}
  \resizebox{\linewidth}{!}{\input{tables/#1}}
  \end{minipage}\par\vspace{10pt}}

\title{Exploratory Data Analysis:\\ Remote Work and Urban Housing Prices in Eight Sun Belt Metros}
\author[1]{Sai Krishna Varshith Danda}
\author[1]{Gayathri Kota}
\author[1]{Eshitha Chowdary Kaka}
\author[1]{Zain Khan}
\affil[1]{Arizona State University, Tempe, AZ 85281, USA}
\date{October 11, 2026}

\begin{document}
\maketitle

\paragraph{Summary.} We ask whether the flattening of the within-metro house-price gradient seen in 2020--22 persisted
through 2025 in Phoenix, Austin, Atlanta, Dallas, Tampa, Charlotte, Nashville and Las Vegas, using ZIP-level Zillow prices
and rents joined to Census data. This exploratory analysis describes the data and does not test our hypotheses formally.
The hypotheses are not holding uniformly. A flattening in 2020--22 that then persists is clear only in Atlanta; Austin's growth
slope turns positive only in 2023--25, Tampa's flattening does not hold, Charlotte, Dallas and Las Vegas show no flattening,
Phoenix's negative gradient depends on far-out ZIP codes, and Nashville's outer ZIP codes were already outgrowing the center
before 2020. Pre-pandemic work-from-home share overlaps strongly with income, and more expensive ZIP codes tended to grow more
slowly before 2020. Within metros, the relationship between work-from-home share and later growth falls to about zero ($+$0.03
in 2020--22 and $+$0.13 in 2023--25, intervals including zero) once pre-2020 growth is controlled, so the apparent post-2020
relationship mostly reflects the earlier trend. In five of the eight metros the highest-value ZIP codes are concentrated on one
side of the center, and a single distance variable fits poorly in some of them. Data problems cluster at the metro edge.

\section{Data Sources, Preparation, and Cleaning}

\subsection{Sources}
Our outcome is the Zillow Home Value Index (ZHVI) \cite{zhvi}: the typical value of homes in the 35th to 65th percentile for
a ZIP code, all homes (single-family and condominium), smoothed and seasonally adjusted, monthly from January 2000 to August
2026. The file covers 26,268 ZIP codes. Because ZHVI comes from Zillow's valuation models and is smoothed, it blurs the
month-to-month timing of changes; we treat that as a limitation. Our secondary outcome is the Zillow Observed Rent Index
(ZORI) \cite{zori}, all homes plus multifamily, smoothed, monthly from January 2015, for 8,459 ZIP codes nationally. Rents do
not depend directly on mortgage financing, so we use them as a check on whether post-2022 price patterns reflect interest
rates or location preferences.

Neighborhood characteristics come from the American Community Survey (ACS) 5-year estimates \cite{acs19,acs23} through the
Census Data API, at ZIP Code Tabulation Area (ZCTA) level. We pull tables B08006 (means of transportation to work, including
``worked from home''), C24010 (occupation by sex), B19013 (median household income), B25024 (units in structure) and B01003
(total population), with margins of error, for two releases: 2015--2019, our pre-pandemic baseline, and 2019--2023, used
only to describe current composition because it mixes behavior from before and after 2020. Distances use ZCTA internal points
from the Census Gazetteer \cite{gaz}, a county-based metro delineation \cite{cbsa} to cross-check metro membership, and
cartographic ZCTA boundaries \cite{zcta} for maps. City-hall coordinates, the centers from which we measure distance, were
compiled by hand from \cite{wiki,osm} (Table~\ref{tab:centers}).

\subsection{Preparation and cleaning}
\textbf{Selecting the metros.} The ZHVI file labels each ZIP code with a metro, but the label is blank for 4,712 ZIP codes
and several names collide (Austin, Minnesota; Las Vegas, New Mexico; Charlottesville, Virginia). We kept the ZIP codes whose
label is one of our eight metros (1,099 ZIP codes) and checked the choice against the Census county-based delineation. The
two methods agree on 1,087 ZIP codes. Twelve ZIP codes are in the Zillow metro but in a county outside the 2023 delineation
(for example Marble Falls, Mabank and Perrin, Texas), and five are in a delineation county but have a blank label (four in
Hickman County, Tennessee and one in Lumpkin County, Georgia). The disagreements look like a vintage difference between
delineations: the Census file is dated July 2023 and uses different titles for Atlanta, Austin and Las Vegas than Zillow
does, so those three had to be matched by CBSA code. We kept Zillow's label, flagged the 12 ZIP codes for a sensitivity
check, and left the 5 out.

\textbf{Reshaping and missing data.} The wide ZHVI file (one column per month) was reshaped to a ZIP-by-month panel of
351,680 rows (1,099 ZIP codes by 320 months), keeping every month. The 2018--2025 window is applied at analysis time. We do
not drop ZIP codes with gaps by default. Of the 1,099 ZIP codes, 1,072 (97.5\%) have a complete series for 2018--2025; 17
start after January 2018 and 10 have gaps in the middle (September 2020 to September 2021). Table~\ref{tab:coverage}
shows coverage by metro. ZIP codes are kept as five-character text so that leading zeros are not lost.

\textbf{Growth measures.} Growth is the year-over-year change, with December 2019 as the reference month (price index
$=100$ in December 2019). We cut the window into three periods: 2018--2019, 2020--2022 and 2023--2025, and, because the
first period is two years and the others three, we annualize growth in all figures so the periods compare. All plots use the
same 1,082 ZIP codes that have a value at every date (``balanced sample''), so a change in a slope between two panels cannot
come from a change in who is in the sample. Prices are nominal.

\textbf{Distance to the center.} Our primary distance is from each ZIP code's ZCTA internal point to the city hall of the
metro's first-named city. We compiled the eight coordinate pairs by hand, with a source recorded for each
(Table~\ref{tab:centers}), instead of geocoding. Dallas--Fort Worth and Tampa--St.\ Petersburg are two-center metros, so for
those two we also add Fort Worth and St.\ Petersburg and compute distance to the nearer center. A third definition, the
centroid of the ZCTA containing the city-hall address, is kept as a robustness check; it sits within 1.5 miles of city hall in
every metro, so it changes almost nothing (rank correlation 0.986 or higher with city-hall distance). One check failed:
Dallas City Hall's coordinate falls in ZCTA 75202, while its address ZIP is 75201 (adjacent downtown ZIP codes); the effect is
under a mile and we left it.

\textbf{ZIP to ZCTA joins and reliability.} ZIP codes join to ZCTAs by identical code. In the 2015--2019 release 1,092 of 1,099
(99.4\%) match; the seven that do not (in Frisco, McKinney, Arlington, Surprise and Las Vegas, plus Mesquite, Nevada) are
mostly fast-growing suburbs that were probably created after the 2010 boundaries. All 1,099 match in the 2019--2023 release.
The ACS work-from-home share is noisy in small ZCTAs, so we compute a margin of error for the share and flag it
\emph{unreliable} when the margin exceeds half the estimate (a coefficient of variation above 0.30). That flags 248 of 1,092
ZIP codes (23\%); they sit a median 39 miles from the center against 17 miles for the rest. A rule based on the number of
workers would have flagged only 4. Census sentinel codes for suppressed values were set to missing.

\textbf{Rents.} ZORI exists for only 799 of our 1,099 ZIP codes (72.7\%), and only 398 have every month of 2018--2025; the
set also grows over time (for example Dallas has rents for 78 ZIP codes in January 2018 and 196 in December 2025). Rent
trends are therefore computed on a fixed set of ZIP codes (those with rents in both December 2019 and December 2025).
Section~\ref{sec:robust} explains why this makes rents a check for the inner and middle metro only.

\tab{t_coverage.tex}{tab:coverage}{Coverage and data quality by metro. ``Complete ZHVI'' = every month 2018--2025;
``Any ZORI'' = at least one rent observation; ``WFH share unreliable'' = margin of error above half the 2015--2019
work-from-home share; ``Metro column only'' = Zillow assigns the ZIP code to the metro but the county-based delineation does
not.}

\tab{t_centers.tex}{tab:centers}{City-hall coordinates, compiled by hand with a source for each (retrieved October 6,
2026). Fort Worth and St.\ Petersburg are the second centers for the two-center metros. An OpenStreetMap address lookup
agrees with the Wikipedia point to within 0.25 miles where both exist.}

\subsection{What we found while cleaning: problems sit at the edge}
Most of the data-quality problems we found are more common far from the center (Figure~\ref{fig:flags}). Beyond 30 miles,
64\% of ZIP codes have no rent series (6\% within 10 miles), 53\% have an unreliable work-from-home share (11\% within 10
miles), 6\% have an incomplete price series and 4\% are in Zillow's metro but not the county delineation (none within 20
miles). Two flags do not follow the pattern. The seven ZIP codes with no 2015--2019 ACS row are not tied to distance
(a within-metro rank test gives $p=0.39$, and only one of the seven is 30 or more miles out); with seven ZIP codes the test has
little power, so we say there is no evidence of a peripheral pattern for them, not that there is none. Patchy ZORI is
the second exception, for a mechanical reason: outer ZIP codes mostly have no rents at all, so only inner ones can be
patchy. The practical consequence is that dropping flagged ZIP codes trims the outer edge of every metro, which is where a
price-distance gradient is estimated. We therefore keep every ZIP code by default, carry the flags in the data, and show
results for the full and the flag-free samples side by side (Section~\ref{sec:robust}).

\fig{2_1c_flags_by_distance.jpg}{fig:flags}{Share of ZIP codes in each distance band carrying each data-quality flag. Darker
means more ZIP codes. Four of the six rows rise with distance; missing 2015--2019 ACS rows are few and not tied to
distance, and patchy ZORI falls with distance because outer ZIP codes mostly have no ZORI at all.}

\subsection{Reproducibility}
All of the above is produced by eight scripts run in order by one command (about 20 seconds once the raw files are in
place). Each writes a data-quality report that this section draws on. The Census API key is read from an environment
variable and never stored. Every file used is listed with its URL, vintage, size and retrieval date in the repository's
\texttt{docs/SOURCES.md}.

\section{Exploratory Data Analysis}

\subsection{Distributions}
Home values are strongly right-skewed (Figure~\ref{fig:logcheck}): skewness is about 2.6 in dollars and 0.4 in logs, so we
plot prices on a log scale throughout. Median ZIP values in December 2019 run from about \$204k (Charlotte) to \$331k
(Austin). Every metro jumped by December 2022 (Austin's median reached \$525k, Phoenix's \$460k) and then diverged: Austin's
fell to \$446k by December 2025 and Phoenix's to \$441k, while Atlanta, Charlotte and Nashville edged up
(Figure~\ref{fig:dist}). Annualized growth in 2020--22 had metro medians of 12.6\% (Las Vegas) to 17.5\% (Tampa) a year; in
2023--25 they ranged from $-$5.1\% (Austin) to $+$1.5\% (Charlotte). Inside a metro, the middle half of ZIP codes spans a
range of 1.4 to 1.8 times in value, and the spread of growth across ZIP codes (a middle-half range of roughly 1 to 5
points a year) is as large as many of the differences between metro medians, so a metro average hides a lot. These are nominal
values from a smoothed, model-based index, so timing is blurred and inflation is not removed.

\fig{2_1a_distributions.jpg}{fig:dist}{Home values (top, log scale) at three dates and annualized growth (bottom) in three
periods, by metro. Each box is the middle half of the metro's ZIP codes, the white line is the median, and whiskers extend
1.5 times the middle-half range.}

\fig{2_1b_log_check.jpg}{fig:logcheck}{Why we use a log scale: the distribution of December 2019 home values on a dollar
axis (left) and a log axis (right).}

\subsection{The gradient (sub-question 1)}
\label{sec:gradient}
The monocentric model predicts that home values fall with distance from the center and that the gradient flattens when
commuting gets cheaper. For each metro we fit home value (in logs) and annualized growth against distance from city hall.
Three cautions apply to everything in this subsection. First, the intervals are plain ordinary least squares and ignore
that nearby ZIP codes are correlated, so they are probably too narrow. Second, these are simple fits of one outcome on
distance, without controls or metro-by-month effects; they are not the planned event study and say nothing about what
remote work did. Third, the smooth (LOWESS) curves are drawn only over the range holding 95\% of a metro's ZIP codes and rest
on few ZIP codes beyond roughly 40 to 60 miles, so their far ends should not be over-read.

Home values fall with distance in six of the eight metros, but the strength varies (Table~\ref{tab:slopes},
Figure~\ref{fig:slopes}). Charlotte (about $-$18\% per 10 miles, R$^2$ about 0.5), Austin ($-$15\%) and Nashville ($-$13\%)
have steep, fairly tight gradients. Atlanta, Tampa and Phoenix are weaker ($-$4\% to $-$9\%, R$^2$ under 0.12). Dallas and Las
Vegas have no linear gradient at all (slopes near zero, R$^2$ below 0.01). The prediction was a flattening in 2020--22 that
then holds. Only Atlanta follows it and keeps going: its level slope moves from $-$7.7\% to $-$6.0\% to $-$3.0\% and its
growth slope from $-$0.44 to $+$0.69 to $+$1.03 points per year per 10 miles (Figure~\ref{fig:atlanta}). The growth-slope change
from before 2020 to 2020--22 is clear (the intervals do not overlap); the level-slope intervals at the three dates overlap, and
we did not test the change. Austin's growth slope is positive only in 2023--25 ($+$0.63), though it is not clearly different from
its 2020--22 value ($+$0.21, interval $-$0.22 to $+$0.64). Tampa's growth slope is $+$0.57 in 2020--22 and $-$0.32 in 2023--25
(interval $-$0.68 to $+$0.03), so its flattening did not hold. Phoenix goes the other way on growth: its growth slope falls from
$+$0.46 to $-$0.30. Its level slope is $-$4.3\% at the first date and $-$5.6\% at the last, with overlapping intervals, and the
gradient disappears when the 16 ZIP codes with unreliable exposure estimates (median distance 56 miles) are dropped
(Section~\ref{sec:robust}; Figure~\ref{fig:phoenix}). Charlotte shows no detectable change, and Nashville's outer ZIP codes were
already outgrowing the center before 2020 ($+$1.02). We show Atlanta and Phoenix in detail because one follows the prediction
and the other reverses it; Figure~\ref{fig:slopes} summarizes all eight metros, and the per-metro figures for the other six
metros and the Tampa two-center growth figure are in the project repository (\texttt{outputs/figures/}).

The surprise is how little of this is one story. In Phoenix and Las Vegas the smooth curves are humped, with values peaking about
24 and 15 miles from city hall and falling beyond (Figure~\ref{fig:phoenix} for Phoenix); a quadratic term is significant in both
($p=0.004$ and $p<0.001$ for December 2025), and a straight line flattens the hump into a near-zero or mildly negative slope.
The Phoenix hump rests on sparse far-out ZIP codes, though: dropping the 20 ZIP codes beyond 40 miles gives a level slope of
$+$8.1\% per 10 miles (interval $+$1.4 to $+$15.1), so inside 40 miles values rise with distance. Level slopes at different dates
have overlapping intervals in every metro, and we did not test whether a slope changed between dates. By a rough rule (the
largest shift between periods exceeds the width of one interval; a comparison, not a test), the growth slopes of Atlanta,
Nashville, Tampa, Phoenix and, marginally, Austin shifted; those of Charlotte, Dallas and Las Vegas did not. Level slopes move
little between December 2019 and December 2025 (within about 1.5 points everywhere except Atlanta, $+$4.7, and Austin,
$+$2.2), which is arithmetic: a growth slope of $+$1 point per year per 10 miles adds only about 3 points to the level slope
over three years. We also cannot compare with dense metros, because none are in our sample, so the plan's expectation that
flattening is smaller here is untested.

\tab{t_slopes.tex}{tab:slopes}{Distance gradients by metro (OLS, 95\% intervals in brackets): home value (log) at three dates, as
\% change per 10 miles from the center; annualized growth in three periods, as points per year per 10 miles. The rows marked
``nearest of 2'' use distance to the nearer of the two centers.}

\fig{2_2_slope_summary.jpg}{fig:slopes}{The slopes of Table~\ref{tab:slopes} with 95\% intervals. Negative means values or growth fall with
distance. A monocentric post-2020 flattening would move the level slope toward zero and turn the growth slope positive.
Intervals that cross zero cannot be told from a flat gradient.}

\fig{2_2_gradient_atlanta.jpg}{fig:atlanta}{Atlanta: home value (log scale, top) and annualized growth (bottom) against
distance from city hall. Each dot is a ZIP code; the solid line is a linear fit and the dashed line a smooth fit.}

\fig{2_2_gradient_phoenix.jpg}{fig:phoenix}{Phoenix, same layout. The smooth fit peaks about 24 miles from city hall; triangles
mark two ZIP codes beyond the axis range.}

\paragraph{Does a second center change the picture?} For Dallas--Fort Worth and Tampa--St.\ Petersburg we recompute distance
to the nearer of the two city halls. This does not recover a monocentric gradient (Figures~\ref{fig:dalv} and~\ref{fig:dalg}).
For Dallas, R$^2$ for home values stays at about 0.005 either way and the slope changes sign ($-$1.9\% to $+$1.5\% per 10
miles) with both intervals spanning zero. For Tampa, R$^2$ rises a little (0.065 to 0.080; the level R$^2$ values are averaged
over the three dates). For growth the nearest-center version helps most in 2020--22 (R$^2$ 0.004 to 0.019 in Dallas, 0.117 to
0.163 in Tampa). Fitting each center's ZIP codes separately points to why a pooled line fails in Dallas--Fort Worth: Fort Worth
ZIP codes rise in value with distance from Fort Worth's city hall by about $+$8\% to $+$10\% per 10 miles (R$^2$ about 0.10)
while Dallas ZIP codes are flat or slightly negative. In Dallas--Fort Worth the two cores sit at different price levels (median
December 2019 value of \$221k within 10 miles of Dallas City Hall and \$183k around Fort Worth's); in Tampa--St.\ Petersburg
they do not (\$220k and \$227k), and the St.\ Petersburg group is too small to say anything (39 ZIP codes; its level-slope
interval runs from about $-$15\% to $+$34\% per 10 miles). The per-center growth slopes in Dallas--Fort Worth also move in
opposite directions (Dallas from $+$0.47 in 2020--22 to $-$0.25 in 2023--25, Fort Worth from $-$0.29 to $+$0.37), an exploratory
observation from small groups. This shows that a second center matters for the price level in Dallas--Fort Worth, not that two
city halls explain the gradient; in Tampa the second center changes little. A city hall is only a proxy for the center of
employment, and we have no job-location data, so suburban job centers are not in the test.

\fig{2_2_two_center_level_dallas.jpg}{fig:dalv}{Dallas--Fort Worth, home values. Row A: distance to Dallas City Hall. Row B:
distance to the nearer of the two city halls, one pooled fit. Row C: the same distances with a separate fit for each center's
ZIP codes (line color matches the dots). Dot color is the center each ZIP code is nearer to, in every row.}

\fig{2_2_two_center_growth_dallas.jpg}{fig:dalg}{Dallas--Fort Worth, annualized growth, same layout as
Figure~\ref{fig:dalv}.}

\subsection{Spatial structure}
The maps ask whether expensive ZIP codes form rings around the center, as the monocentric model predicts, or separate
patches, as a polycentric metro would. Figure~\ref{fig:lvlmap} shows December 2019 home value in fifths within each metro. We
also tested whether each metro's highest-value fifth is concentrated in one direction from the center, by comparing its
compass-bearing concentration with that of random fifths of the same ZIP codes (5,000 permutations; Table~\ref{tab:wedge}). In
five of the eight metros the highest-value ZIP codes are concentrated on one side of the center ($p<0.01$ in each): northeast of
Phoenix's city hall, northwest of Austin's, north of Atlanta's, northeast of the two Dallas--Fort Worth halls, and south of
Nashville's. In Dallas--Fort Worth both cores sit in lower fifths (63\% of the ZIP codes within 10 miles of either hall are in
the bottom two fifths, against 40\% expected without a pattern). Tampa, Charlotte and Las Vegas show no clear direction
($p=0.77$, 0.15 and 0.08); in Las Vegas the dark area to the northwest on the map is a few very large ZIP codes, and the
highest-value fifth points southwest. A wedge is not the shape the monocentric model predicts, which would help explain why one
distance variable fits poorly where the wedge is strongest (Phoenix, Dallas); but Austin and Nashville have wedges and still fit
moderately well (R$^2$ 0.15 to 0.24 and 0.26 to 0.37 in Section~\ref{sec:gradient}), so the wedge alone does not account for the
fit.

The growth maps (Figures~\ref{fig:mapatl}--\ref{fig:mapdal} show three; the maps for the other five metros are in the project
repository and notebook) show growth relative to each metro's own median, on one color scale for all eight metros. In Atlanta,
Austin and Nashville the ZIP codes within 10 miles of the center grew more slowly than the metro in 2023--25 and those beyond 30
miles faster (median differences of $-$2.9 against $+$1.5, $-$1.3 against $+$0.4 and $-$1.7 against $+$0.8 points a year; $p<0.01$
in each, Mann-Whitney), which matches their positive 2023--25 slopes. Charlotte's core is only mildly slow ($-$0.7, and its
difference from the outer ZIP codes is not significant), and its far-south ZIP codes are also slow ($-$0.7 against $+$0.6
elsewhere beyond 30 miles, 10 ZIP codes), which may be why its slope is near zero. Before 2020 the pattern was different:
Atlanta's ZIP codes north of city hall were slow ($-$3.7 points a year against the metro) and those to the south and southeast
fast ($+$1.1 and $+$2.2); the Dallas and Fort Worth cores grew faster than their surroundings ($+$2.4 within 10 miles of either
hall against $-$0.5 beyond), with the northern suburbs slower ($-$3.5 north against $+$1.4 south of the midpoint, beyond 10
miles); and Phoenix's far southeast was fastest ($+$2.6), while its inner core was only modestly above the metro median
($+$0.9). In seven of eight metros (all but Austin) growth varied more across ZIP codes before 2020 than in 2020--22 (the standard
deviation of annualized growth was 4.9 against 3.5 points a year in Atlanta and 4.2 against 1.9 in Dallas), so the post-2020 maps
are, if anything, quieter than the pre-2020 ones; the pre-2020 window is two years long, though, which by itself inflates
annualized dispersion. The maps are descriptive: growth is relative to the metro median so metro-wide movements are removed;
large, thinly populated ZIP codes dominate the area (Las Vegas is hard to read for this reason); and we have not measured spatial
clustering or explained why the wedges form.

\tab{t_wedge.tex}{tab:wedge}{Is the highest-value fifth (December 2019) concentrated in one direction from the center? Concentration R is 0 if
the ZIP codes point in all directions equally and 1 if they all point one way; the permutation p-value compares R with random
fifths of the same ZIP codes (5,000 draws). Bearings use the ZCTA internal point; for Dallas--Fort Worth the origin is the
midpoint of the two city halls. The Las Vegas direction is southwest, not northwest.}

\fig{2_3b_level_quintile_maps.jpg}{fig:lvlmap}{December 2019 home value as fifths within each metro. Red star = city hall (two
in Dallas--Fort Worth and Tampa--St.\ Petersburg). Fifths are within each metro, so colors compare places inside a metro, not
prices across metros.}

\fig{2_3a_growth_map_atlanta.jpg}{fig:mapatl}{Atlanta: annualized growth relative to the metro median, by period. Warm =
faster than the metro, cool = slower. The same color scale is used for all metros.}

\fig{2_3a_growth_map_phoenix.jpg}{fig:mapphx}{Phoenix, same layout.}

\fig{2_3a_growth_map_dallas.jpg}{fig:mapdal}{Dallas--Fort Worth, same layout; the two stars are the two city halls.}

\subsection{Pre-pandemic work-from-home share and growth (sub-question 2)}
The plan predicted a positive relationship between a ZIP code's pre-2020 work-from-home potential and its price growth
afterward, strongest in 2021--22 and weakening but not disappearing through 2025. We measure exposure by the share of workers
who worked from home in the 2015--2019 ACS.

The first thing to know is that high-work-from-home and high-income ZIP codes overlap strongly (Figures~\ref{fig:confound}
and~\ref{fig:corr}). Within a metro, the work-from-home share has a rank correlation of $+$0.66 with income, $+$0.74 with home
value and $+$0.74 with the share of management and professional workers. More expensive ZIP codes tended to grow more slowly
before 2020 (growth against December 2019 value, rank correlation $-$0.72). So the pre-2020 relationship between exposure and
growth is clearly negative ($-$1.48 points a year per $+$5 percentage points of work-from-home share, R$^2$ 0.13;
Figure~\ref{fig:expo}, Table~\ref{tab:expo}): before remote work could matter, affluent ZIP codes tended to be on a slower path.
Within metros the slope is $-$0.30 in 2020--22 and $+$0.40 in 2023--25 (pooled: $+$0.17 and $-$0.10).

\paragraph{Controlling for the pre-2020 trend.} To see how much of the later slopes is the earlier trend, we regressed 2020--22
and 2023--25 growth on work-from-home share within metro (all variables measured against the metro median), adding pre-2020
growth and then income (1,073 ZIP codes with income data; robust standard errors; Table~\ref{tab:expoctl}). Adding pre-2020
growth reduces the slopes to $+$0.03 (interval $-$0.13 to $+$0.20) for 2020--22 and $+$0.13 ($-$0.03 to $+$0.30) for 2023--25;
adding income as well gives $-$0.15 ($-$0.36 to $+$0.06) and 0.00 ($-$0.15 to $+$0.15). All four intervals include zero. The
later slopes are therefore largely accounted for by the earlier trend, and we find no relationship with exposure that is separate
from it. This is an exploratory regression, not the planned event study, and the intervals ignore spatial correlation.

The slope does rise across the three periods ($-$1.48, $-$0.30, $+$0.40), but its only positive value is in 2023--25, whereas the
plan expected the strongest relationship in 2021--22, and it falls to about zero once pre-2020 growth is controlled. We therefore
do not read it as support for the hypothesis. A slope that is less negative than before could reflect exposure, the same affluent
ZIP codes flattening out, or anything else that hit expensive homes differently, such as the 2022 rate shock; the control result
favors the second.

Two further points. The strongest relationship in the figure is the pre-2020 placebo, and the plan's predicted positive
relationship in 2020--22 does not appear within metros (the pooled version, which includes differences between metros, is
$+$0.17 with an interval from $-$0.01 to $+$0.35). Dropping the 233 ZIP codes in this sample with an unreliable estimate (hollow
diamonds in Figure~\ref{fig:expo}) barely changes the slopes ($-$1.41, $-$0.26, $+$0.52), though those ZIP codes scatter somewhat
more widely (residual spread 8\% to 24\% larger). The controls above handle only pre-2020 growth and income; exposure also tracks
occupation mix (rank correlation $+$0.74), and we have no education measure. A model of the post-2020 relationship should include
pre-period growth and test for pre-trends. The same three cautions apply as in Section~\ref{sec:gradient}.

\fig{2_4c_exposure_confounds.jpg}{fig:confound}{What the work-from-home share is entangled with, within metro: income (left)
and distance (right). Hollow diamonds are ZIP codes whose work-from-home share is unreliable.}

\fig{2_4b_exposure_corr.jpg}{fig:corr}{Within-metro rank correlations among the exposure measure, its likely confounders and
growth. Each variable is first converted to its percentile within its own metro.}

\fig{2_4a_exposure_scatter.jpg}{fig:expo}{Growth against pre-2020 work-from-home share. Top: within metro (both measured against
the metro median). Bottom: metros pooled. Left column is a check on growth before 2020. Dots are ZIP codes (hollow diamonds
unreliable), large dots are means of eighths of the ZIP codes with 95\% intervals, solid line is a linear fit and dashed line a
smooth fit.}

\tab{t_exposure.tex}{tab:expo}{Slopes of annualized growth on work-from-home share, in points per year per $+$5 percentage
points, with 95\% intervals (OLS), all ZIP codes and ZIP codes with a reliable estimate only.}

\tab{t_exposure_controls.tex}{tab:expoctl}{Slope of annualized growth on work-from-home share within metro, in points per year per $+$5 percentage points, with
95\% intervals (robust standard errors), as pre-2020 growth and income are added as controls. All variables are measured
against the metro median; 1,073 ZIP codes.}

\subsection{Price paths and groups of ZIP codes (sub-question 3)}
Every metro followed the same broad shape (Figure~\ref{fig:paths}): a steep rise from 2020 to 2022 (metro median ZIP indexes of
about 143 to 162 by December 2022, December 2019 $=100$) and then a flat or falling stretch, with Austin falling furthest. Within
most metros the distance bands move together: in Phoenix, Dallas and Charlotte each band's median stays within about 4\% of the
metro as a whole throughout (Figure~\ref{fig:relpaths}).

Atlanta is the one clear case of inner ZIP codes losing ground to outer ones (Table~\ref{tab:donut}). Its innermost band (0--10
miles) fell to 85\% of the metro's median index by December 2022 and 83\% by December 2025 (bootstrap 95\% interval 80.5 to 87.2),
while the 30-plus-mile band rose to about 105\% (103.6 to 106.2). That is the pattern the literature calls a donut, but the inner
decline began before 2020: the inner band's index relative to the metro was already falling, from 103.3 in January 2018 to 100 at
December 2019. The post-2020 fall (15 points in three years) is much steeper than the earlier one (3 points in two years), so it
is not a simple continuation, but it is not new either. Nashville and Austin show weaker versions in the same direction.
Nashville's inner band fell from 107.2 in January 2018 to 100 at December 2019, rose to 102.9 by December 2022 and ended at 97.1
(interval 93.6 to 100.0), so its decline predates 2020; Austin's inner band ended at 95.6 (92.3 to 97.4) and its outer band at
101.4 (95.5 to 106.8), an interval that spans 100. Tampa does not fit: its inner and outer bands both gained on the middle bands
(105.8 and 106.7 in December 2025, against 97.7 and 98.8 for the 10--20 and 20--30 mile bands). In Austin and Las Vegas the outer
band swings (Austin's dipped to about 92\% in 2022 and recovered; Las Vegas's, with only 7 ZIP codes, fell to about 94\% in 2022
and ended near 96\%).

Grouping the 1,072 ZIP codes with complete series by the shape of their path (k-means on the log index) gives four groups
(Figure~\ref{fig:clusters}, Table~\ref{tab:clusters}). Group 3 (89 ZIP codes, 74\% of them in Austin) shows a boom and a correction,
ending about 33\% above December 2019 after a peak near 169. Group 4 (48 ZIP codes, median 4 miles from city hall, 18\%
single-family detached homes, mostly Atlanta and Dallas) rose least, peaking near 128. Groups 1 and 2 (520 and 415 ZIP codes)
differ mainly in amplitude. The groups follow metro more than distance (Cram\'er's V is 0.60 for metro and 0.18 for distance
band; median distances are 22, 20, 15 and 4 miles), so only the small urban-core group is distinctively central; housing type
separates only Group 4 (single-family share 18\% against 65\% to 71\% in the other groups). The plan expected three to five groups
including delayed appreciation and stagnation; we find a boom-and-correction group and a low-appreciation group but no
delayed-appreciation group. This comes with a caveat: clustering on the log index is dominated by overall amplitude (the first
principal component alone explains 71\% of the variance in the paths), so shape types such as delayed appreciation may not
separate, and we did not test shape-only clustering.

The groups are only weakly separated. The silhouette width is 0.22 to 0.38 across two to eight groups (0.36 at four, no better
than 0.375 at two), 25\% of ZIP codes have a silhouette below 0.25, and the four-group solution is stable to the starting point
but not always to resampling (adjusted Rand index against the reported solution: median 0.85, fifth percentile 0.45 across 40
bootstrap resamples). We therefore cannot say that four is the right number of groups or that the groups are real types and not a
partition of a continuum; the dominance of one dimension is what a continuum looks like. The 27 ZIP codes with incomplete series,
mostly in the outer ring, are left out, and the band medians in Las Vegas and the outer bands of Austin and Tampa rest on 4 to 15
ZIP codes.

\fig{2_5a_trajectories.jpg}{fig:paths}{Price paths of every ZIP code (thin lines) by distance band, with band medians
(thick), December 2019 $=100$, log scale.}

\fig{2_5b_relative_paths.jpg}{fig:relpaths}{Each ZIP code's index divided by its metro's median index, median by distance band.
If outer ZIP codes gained on inner ones after 2020, the outer-band lines rise above the inner ones.}

\tabsmall{t_donut.tex}{tab:donut}{Median price index of the innermost and outermost distance bands relative to the metro (each ZIP code's index, December
2019 $=100$, divided by its metro's median index, so 100 means the metro as a whole; every band is at 100 in December 2019 by
construction), with a bootstrap 95\% interval for December 2025 (2,000 resamples of the ZIP codes in the band).}

\fig{2_5c_clusters.jpg}{fig:clusters}{k-means grouping of ZIP codes by price path: average silhouette width by number of groups
(left), the mean path of each group with a one-standard-deviation band (middle), and the metro mix of each group (right).}

\tab{t_clusters.tex}{tab:clusters}{What each group looks like (medians): distance from city hall, December 2019 home value, 2015--2019
household income and work-from-home share, single-family detached share of homes, and annualized growth in three periods (\%).}

\subsection{Robustness}
\label{sec:robust}
\textbf{Flagged ZIP codes.} Dropping the flagged ZIP codes (834 of 1,082 remain) changes the picture more for price levels than for
growth (Figure~\ref{fig:flagfree}). Level slopes move by up to about 6 points per 10 miles (flatter on average by about 2 points),
consistent with the flagged ZIP codes being mostly cheaper (median December 2019 value \$213k against \$269k) and farther out
(Section 1.3). Growth slopes change by 0.00 on average (range $-$0.37 to $+$0.59 points per year per 10 miles). Phoenix is the case
that matters: with the flagged ZIP codes removed its level slope goes from $-$4.3\% to $+$0.5\% per 10 miles and all three growth
slopes are near zero ($+$0.09, $+$0.03, $-$0.09). Only one flag does this: dropping the 16 ZIP codes with an unreliable
work-from-home share (median distance 56 miles) gives the $+$0.5\%, while dropping each of the other three flags (one to three ZIP
codes) changes nothing; dropping the 20 ZIP codes beyond 40 miles gives $+$8.1\% (interval $+$1.4 to $+$15.1). So Phoenix's
negative gradient comes from far-out ZIP codes. Atlanta's pattern survives ($+$1.25 in 2020--22 against $+$0.69 with all ZIP
codes; the intervals overlap); Nashville's and Tampa's growth slopes hold up; and Dallas's 2023--25 slope is $-$0.35 (interval
$-$0.57 to $-$0.12) when flagged ZIP codes are dropped, against $-$0.03 with all ZIP codes, so the Dallas result depends on the
sample. The flag-free subset is a selected sample, neither more nor less correct than the full one.

\textbf{Center definition.} Using the ZCTA centroid of the city-hall ZIP code instead of city hall changes no slope by more than
0.5 percentage points (levels) or 0.06 points per year (growth).

\textbf{Rents.} Rents were meant to check whether post-2022 price patterns reflect mortgage rates, but ZORI is weakest where the
gradient question lives. Only 36\% of ZIP codes 30 or more miles from the center have any rent series, against 94\% within 10 miles
(Table~\ref{tab:zori}); only 13\% of the outer ZIP codes have rents at both December 2019 and December 2025. Coverage is also
selective: among ZIP codes 20 or more miles out, those with rents have a median population of 35,786 against 4,872, and 9.0\% against
0.3\% of homes in buildings of five or more units. The rent check is therefore an inner and middle metro check; it can support or
weaken a rate-shock explanation there but cannot rule one out for the outer ring. On the 478 ZIP codes with rents at both ends,
price growth rose with distance in 2020--22 ($+$0.83 points per year per 10 miles) and rents did too ($+$0.41). In 2023--25 price
growth no longer rose with distance ($-$0.07) while rent growth still did ($+$0.33; Figure~\ref{fig:rent}). This is the kind of
contrast the rent check was designed to look for, since rents do not depend directly on mortgage financing. It is not proof:
rents are smoothed and lag prices, so they may still be adjusting to the 2020--22 shift; the sample is selective; and the
comparison uses the same simple distance fits as above.

\tabsmall{t_zori.tex}{tab:zori}{Rent (ZORI) coverage by distance from city hall, as a share of the ZHVI ZIP codes in each band.}

\fig{2_6a_flag_free.jpg}{fig:flagfree}{Gradient slopes on all ZIP codes (blue) and on flag-free ZIP codes only (orange). Dots are
OLS slopes and whiskers 95\% intervals (probably too narrow).}

\fig{2_6c_rent_vs_price.jpg}{fig:rent}{Prices (top, blue) and rents (bottom, orange) on the same 478 ZIP codes, growth minus the
metro median against distance from city hall. Only 41 of these ZIP codes are 30 or more miles out.}

\section{Implications}
\textbf{Does the hypothesis hold?} Not uniformly. The expected flattening of the price gradient in 2020--22 that then persists is
clear only in Atlanta; Austin's growth slope turns positive only in 2023--25, and Tampa's flattening does not hold. Charlotte,
Dallas and Las Vegas show no flattening, Phoenix's negative gradient depends on far-out ZIP codes, and Nashville's outer-ZIP
advantage predates 2020. The prediction for sub-question 2 (more pre-pandemic work-from-home potential, more growth, strongest in
2021--22) is not what we see: within metros the relationship is negative until 2023--25, and it falls to about zero once pre-2020
growth is controlled, so it largely reflects an existing pre-2020 trend that is correlated with income. Sub-question 3's groups
are only weakly separated, partly track metro, and are dominated by one dimension of overall amplitude.

\textbf{What this means for the modeling.} A single linear distance term describes these metros poorly: expensive ZIP codes are
concentrated on one side of the center in five of the eight metros, values hump in Phoenix and Las Vegas, the two cores of
Dallas--Fort Worth sit at different price levels, and Dallas and Las Vegas show no linear gradient. The models should therefore
allow a smooth distance effect and center-specific terms, and the job-center data we set aside is worth bringing back. Because
pre-2020 growth is already clearly related to work-from-home share, and controlling for it removes the later relationship, the
event study should include pre-period growth and income and report the pre-2020 coefficients. Results that depend on the outer ZIP
codes (Phoenix, Dallas) should be shown with and without the flagged ZIP codes, with standard errors clustered for spatial
correlation, and treated as provisional. Rents can check the rate-shock explanation only for the inner and middle metro, on a fixed
set of ZIP codes, and because ZORI is smoothed a lag between rents and prices cannot be ruled out; the plan's metro-by-month fixed
effects remove metro-wide movements such as the 2022 rate shock.

\textbf{What the data cannot answer.} Whether remote work caused any of this; what happened in the outer ring, where rents and
reliable exposure measures are thin; and whether other centers explain the wedges. \textbf{Next.} Add the job-center data, build
the event-study specification with pre-trend tests and spatially clustered errors, and repeat the grouping within metros so that
group membership is not just metro.

\begin{thebibliography}{9}\setlength{\itemsep}{0pt}
\bibitem{zhvi} Zillow Group. Zillow Home Value Index (ZHVI), ZIP code, all homes, mid-tier, smoothed, seasonally adjusted. Zillow
Research, data through August 2026. \url{https://www.zillow.com/research/data/} (retrieved October 6, 2026).
\bibitem{zori} Zillow Group. Zillow Observed Rent Index (ZORI), ZIP code, all homes plus multifamily, smoothed. Zillow Research,
data through August 2026. \url{https://www.zillow.com/research/data/} (retrieved October 6, 2026).
\bibitem{acs19} U.S. Census Bureau. American Community Survey 5-Year Estimates, 2015--2019, tables B08006, C24010, B19013, B25024,
B01003, ZIP Code Tabulation Areas. Census Data API. \url{https://www.census.gov/data/developers/data-sets/acs-5year.html}
(retrieved October 6, 2026).
\bibitem{acs23} U.S. Census Bureau. American Community Survey 5-Year Estimates, 2019--2023, same tables and geography. Census Data
API (retrieved October 6, 2026).
\bibitem{gaz} U.S. Census Bureau. 2024 Gazetteer Files, ZIP Code Tabulation Areas.
\url{https://www.census.gov/geographies/reference-files/time-series/geo/gazetteer-files.html} (retrieved October 6, 2026).
\bibitem{cbsa} U.S. Census Bureau. Core Based Statistical Areas, Metropolitan Divisions, and Combined Statistical Areas, July 2023
(List 1).
\url{https://www2.census.gov/programs-surveys/metro-micro/geographies/reference-files/2023/delineation-files/list1_2023.xlsx}
(retrieved October 6, 2026).
\bibitem{zcta} U.S. Census Bureau. Cartographic Boundary File, 2020 ZIP Code Tabulation Areas, 1:500,000.
\url{https://www2.census.gov/geo/tiger/GENZ2020/shp/cb_2020_us_zcta520_500k.zip} (retrieved October 6, 2026).
\bibitem{wiki} Wikipedia contributors. City hall articles for Phoenix, Austin, Atlanta, Dallas, Tampa, Nashville (Davidson County
Courthouse) and Las Vegas, for building coordinates. \url{https://en.wikipedia.org} (retrieved October 6, 2026).
\bibitem{osm} OpenStreetMap contributors. Nominatim address search, for the Charlotte, Fort Worth and St.\ Petersburg city-hall
points and for coordinate cross-checks. \url{https://nominatim.openstreetmap.org} (retrieved October 6, 2026).
\end{thebibliography}

\end{document}
